\section{Data appendix}\label{app:data}

\subsection{Sample and variable construction}

I pool the twenty quarterly ENAHO employment modules (module 500) for 2021Q1--2025Q4;
the main analysis uses 2021Q1--2024Q4 and the validation uses 2024Q1--2025Q4. The
working-age sample comprises individuals aged 14 to 65 with non-missing education. All
statistics use the module employment weight (\texttt{fac500a}). Departments are the 25
first-level administrative units derived from the first two digits of \texttt{ubigeo};
the urban indicator is derived from the sampling stratum.

\paragraph{Education and skill groups.} Educational attainment is ENAHO variable
\texttt{p301a}. I verified its coding against the weighted attainment distribution.
Low-skilled workers are those with \texttt{p301a} in 1--6 (no education through complete
secondary); high-skilled workers are those with \texttt{p301a} in 7--11 (non-university
higher through postgraduate). Special education (code 12) and missing values are excluded.

\paragraph{Labor income and wages.} For dependent employees, monthly labor income in the
main occupation is the reported gross pay including overtime (\texttt{p524a1}); for
own-account workers it is net profit (\texttt{p530a}). I deflate by the Lima consumer
price index to December 2021 soles. The real hourly wage divides monthly labor income by
$4.345$ times weekly hours in the main occupation (\texttt{p513t}, imputed where missing).
Wages are winsorized at the 0.5th and 99.5th percentiles. As a robustness check I also
construct the imputed, deflated, annualized labor-income aggregate used by
\citet{jimenez_2023}, namely the sum of \texttt{i524a1}, \texttt{d529t}, \texttt{i530a},
and \texttt{d536} divided by twelve; results are unchanged.

\paragraph{Informality.} INEI's official informal-employment indicator (\texttt{ocupinf},
coded 1 for informal and 2 for formal in these files) is present for 2021--2023 but absent
for 2024--2025. I reconstruct a consistent indicator from primitives available in every
year. A wage employee or domestic worker is classified as formal if affiliated to a
pension system (\texttt{p558a1}--\texttt{p558a4}) or holding a formal labor contract
(\texttt{p511a} in 1--6); an employer or own-account worker is classified as formal if the
productive unit is registered (\texttt{p510a1}); unpaid family workers are always informal.
Validated against \texttt{ocupinf} on the overlapping years 2021--2023, this rule agrees
with the official classification for 88.3 percent of employed workers and reproduces the
official informality rate to within about two and a half percentage points. I use the
official indicator wherever available and the reconstruction only for 2024--2025.

\paragraph{Minimum wage and bite.} The statutory minimum wage schedule is taken from the
decrees and verified against the Banco Central de Reserva del Perú series (PN02124PM).
The department Kaitz index is the ratio of the post-reform minimum (1,025 soles) to the
department's pre-reform (2021Q1--2022Q1) weighted median monthly wage of dependent
employees. The fraction-affected measure is the pre-reform weighted share of dependent
employees earning below 1,025 soles.

\subsection{Supplementary results}

Table~\ref{tab:placebo} reports the in-time placebo reforms and confirms that the wage
placebos are null while the employment placebos are contaminated by the pandemic pre-trend.
The doubly robust estimates of \citet{santanna_zhao_2020} appear alongside the two-way
fixed-effects estimates in Table~\ref{tab:main}. Figure~\ref{fig:val2025} shows the
event study around the January-2025 reform used in the out-of-sample validation.

\begin{figure}[H]
\centering
\includegraphics[width=0.9\textwidth]{fig9_validation2025.pdf}
\caption{Event study around the January-2025 reform (out-of-sample validation). Coefficients
on the interaction of the low-skill indicator with quarter indicators, reference quarter
2024Q4. Bands are ninety-five percent confidence intervals clustered by department.}
\label{fig:val2025}
\end{figure}

\begin{table}[H]\centering
\caption{In-time placebo reforms (pre-reform window only)}\label{tab:placebo}
\begin{threeparttable}
\begin{tabular}{lccc}
\toprule
Outcome & Placebo 2021Q2 & Placebo 2021Q3 & Placebo 2021Q4 \\
\midrule
Log real hourly wage & 0.020 & -0.018 & -0.005 \\
 & (0.037) & (0.032) & (0.021) \\
Employment & -0.057$^{***}$ & -0.034$^{***}$ & -0.026$^{***}$ \\
 & (0.016) & (0.007) & (0.008) \\
Formal employment & -0.021$^{*}$ & -0.012 & -0.005 \\
 & (0.011) & (0.013) & (0.007) \\
Self-employment & 0.022$^{*}$ & -0.009 & 0.007 \\
 & (0.012) & (0.007) & (0.009) \\
\bottomrule
\end{tabular}

\begin{tablenotes}\footnotesize
\item Notes: Each entry is the $\mathrm{Low}\times\mathrm{Post}$ coefficient from a placebo
DiD that assigns a fake reform to the indicated quarter, estimated using only pre-reform
data (2021Q1--2022Q1). Standard errors clustered by department in parentheses; all estimates
are weighted by ENAHO survey weights. $^{*}p<0.1$, $^{**}p<0.05$, $^{***}p<0.01$.
\end{tablenotes}
\end{threeparttable}
\end{table}
